  % Custom margins for this section
\chapter{Fistful}
\section*{Pravidlá}
\noindent High Noon a Fistful sú dve rozšírenia BANG! s rovnakou mechanikou udalostných kariet, líšia sa len obsahom kariet.

\textbf{Príprava:} Šerif vyberie jedno rozšírenie. Titulnú kartu (High Noon alebo Fistful) odloží bokom, zvyšné karty zamiešá lícom nadol do samostatného balíčka. Titulnú kartu dá na vrch, balíček otočí lícom hore a umiestni vedľa seba — titulná karta je posledná v balíčku.

\textbf{Priebeh hry:} Od druhého ťahu šerifa: pred každým ťahom odoberie vrchnú kartu balíčka, prečíta efekt nahlas a uloží do stredu stola. Karta na vrchu tejto hromádky je aktívna a efekt platí, kým ju neprekryje ďalšia udalosť. Titulná karta zostáva aktívna do konca hry.

\textbf{Kombinácia oboch rozšírení:} Šerif odloží titulné karty oboch rozšírení, zvyšné zamiešá dokopy, náhodne vyberie 12 (viac pre dlhšiu hru), potom náhodne zvolí jednu titulnú kartu a dá ju na vrch. Balíček otočí lícom hore.

\textbf{Poznámky:} Formulácia „na začiatku ťahu" znamená pred akoukoľvek inou akciou hráča. Pri kartách Odraz a Ruská ruleta možno použiť ľubovoľnú kartu so symbolom Vedle!.
\begin{figure}[!h]
    \centering
    \begin{minipage}[t]{\minipageSize}\centering
      \includegraphics[width=\cardSize\linewidth]{fistful/05_agguato.png}
      \caption[Agguato]{Vzdialenosť medzi akýmikoľvek dvoma hráčmi je 1. Túto vzdialenosť menia iba karty v hre.}
      \altnames{Ambush}
      \skname{Léčka}
    \end{minipage}
    \hfill
     \begin{minipage}[t]{\minipageSize}\centering
      \includegraphics[width=\cardSize\linewidth]{fistful/05_perunpugnodicarte.png}
      \caption[Per un Pugno di Carte]{Na začiatku ťahu je na každého hráča vystrelených toľko BANG!-ov, koľko kariet drží v ruke.}
      \skname{Fistful}
      \altnames{Fistful of Cards}
    \end{minipage}
    \hfill
     \begin{minipage}[t]{\minipageSize}\centering
      \includegraphics[width=\cardSize\linewidth]{fistful/05_ranch.png}
      \caption[Ranch]{Na konci fázy 1 môže hráč raz odhodiť ľubovoľný počet kariet z ruky a dobrať si za ne rovnaký počet nových.}
      \skname{Ranč}
    \end{minipage}
    \hfill
    \begin{minipage}[t]{\minipageSize}\centering
      \includegraphics[width=\cardSize\linewidth]{fistful/05_fratellidisangue.png}
      \caption[Fratelli di Sangue]{Na začiatku ťahu môže hráč darovať jeden svoj život (okrem posledného) inému hráčovi podľa vlastného výberu.}
      \skname{Pokrvní bratia}
      \altnames{Blood Brothers, Pokrvní Bratři}
    \end{minipage}
\end{figure}

\begin{figure}[!h]
    \centering
    \begin{minipage}[t]{\minipageSize}\centering
      \includegraphics[width=\cardSize\linewidth]{fistful/05_ilgiudice.png}
      \caption[Il Giudice]{Nemôžete vyložiť karty pred seba alebo pred iných hráčov.}
      \skname{Sudca}
      \altnames{The Judge, Soudce}
    \end{minipage}
    \hfill
    \begin{minipage}[t]{\minipageSize}\centering
      \includegraphics[width=\cardSize\linewidth]{fistful/05_lazo.png}
      \caption[Lazo]{Karty, ktoré sú vyložené pred hráčmi, neúčinkujú. Platí pre modré aj zelené karty.}
      \skname{Laso}
    \end{minipage}
    \hfill
    \begin{minipage}[t]{\minipageSize}\centering
      \includegraphics[width=\cardSize\linewidth]{fistful/05_leggedelwest.png}
      \caption[Legge del West]{Počas fázy 1 ukáže každý hráč druhú kartu, ktorú si doberá: ak môže, musí ju počas fázy 2 zahrať.}
      \skname{Právo Západu}
      \altnames{Law of the West}
    \end{minipage}
    \hfill
    \begin{minipage}[t]{\minipageSize}\centering
      \includegraphics[width=\cardSize\linewidth]{fistful/05_liquoreforte.png}
      \caption[Liquore Forte]{Každý hráč môže preskočiť fázu 1 a získať tým 1 život.}
      \skname{Pálenka}
      \altnames{Hard Liquor}
    \end{minipage}
\end{figure}

\begin{figure}[!h]
    \centering
    \begin{minipage}[t]{\minipageSize}\centering
      \includegraphics[width=\cardSize\linewidth]{fistful/05_minieraabbandonata.png}
      \caption[Miniera Abbandonata]{Počas fázy 1 si každý hráč berie karty z odhadzovacieho balíčka; ak je prázdny, berie ich z dobracieho balíčka. Počas fázy 3 sa karty odhadzujú lícom nadol na spodok dobracieho balíčka. Takto sa ukladajú iba karty, ktoré sa naozaj odhadzujú, napríklad vo fáze 3 alebo pri Dueli. Efekt vždy platí iba pre hráča, ktorý je práve na ťahu. Napríklad hráč, ktorý sa vyhol efektu Bang!, odhodí kartu Vedle! do bežnej odhadzovacej kopy. Efekt sa nevzťahuje na schopnosti postáv, napr. Jose Delgado. \cite{Fistful}}
      \skname{Opustená baňa}
      \altnames{Abandoned Mine, Opuštěný dúl}
    \end{minipage}
    \hfill
     \begin{minipage}[t]{\minipageSize}\centering
      \includegraphics[width=\cardSize\linewidth]{fistful/05_cecchino.png}
      \caption[Cecchino]{Počas svojho ťahu môže hráč proti inému hráčovi odhodiť dve karty Bang! naraz. Stále ide len o jeden efekt Bang!, teda o stratu 1 života, ale odvrátiť ho možno iba dvoma efektmi Vedle!. Využitie tohto efektu sa nepočíta do limitu Bang!. \cite{Fistful}}
      \skname{Ostreľovač}
      \altnames{Sniper}
    \end{minipage}
    \hfill
    \begin{minipage}[t]{\minipageSize}\centering
      \includegraphics[width=\cardSize\linewidth]{fistful/05_peyote.png}
      \caption[Peyote]{Namiesto ťahania kariet vo fáze 1 tipuje každý hráč, či je karta na vrchu balíčka červená alebo čierna. Ak uhádne, kartu si nechá a môže tipovať znova, inak pokračuje do fázy 2. Hráč si nenecháva ani neuhádnutú kartu. Teda je možné že vo fáze 1 nedostane žiadnu kartu. \cite{Fistful}}
    \end{minipage}
    \hfill
    \begin{minipage}[t]{\minipageSize}\centering
      \includegraphics[width=\cardSize\linewidth]{fistful/05_deadman.png}
      \caption[Dead Man]{Hráč, ktorý vypadol ako prvý, sa počas svojho ťahu vracia do hry s 2 životmi a 2 kartami. Hráč si následne ťahá aj za fázu 1. Do hry sa môžte vrátiť aj keď ste zomreli pred vaším ťahom a Dead Man je práve v hre.}
      \skname{Mŕtvy muž}
      \altnames{Dead Man}
    \end{minipage}
\end{figure}

\begin{figure}[!h]
    \centering
    \begin{minipage}[t]{\minipageSize}\centering
      \includegraphics[width=\cardSize\linewidth]{fistful/05_rimbalzo.png}
      \caption[Rimbalzo]{Každý hráč môže odhodiť kartu BANG! proti karte vyloženej pred akýmkoľvek hráčom: ak hráč nepoužije efekt Vedľa!, zasiahnutá karta sa odhodí. Dostrel nie je potrebný pri využití tohto efektu. \cite{Fistful}}
      \skname{Odrazená Strela}
      \altnames{Ricochet, Odřazená Střela}
    \end{minipage}
    \hfill
    \begin{minipage}[t]{\minipageSize}\centering
      \includegraphics[width=\cardSize\linewidth]{fistful/05_rouletterussa.png}
      \caption[Roulette Russa]{Keď začne účinkovať Ruská ruleta, všetci hráči odhadzujú karty s efektom Vedľa! až kým jeden z nich nemá čo odhodiť: ten stratí 2 životy a ruleta končí.}
      \skname{Ruská ruleta}
      \altnames{Russian Roulette}
    \end{minipage}
    \hfill
    \begin{minipage}[t]{\minipageSize}\centering
      \includegraphics[width=\cardSize\linewidth]{fistful/05_vendetta.png}
      \caption[Vendetta]{Na konci ťahu každý hráč „sejme": ak je karta srdcová, hrá ešte raz (ale už znovu „nesejme").}
      \skname{Vendeta}
    \end{minipage}
       \hfill
    \begin{minipage}[t]{\minipageSize}\centering
    \end{minipage}
\end{figure}
